\documentclass{article}

\usepackage[top=0.5in, bottom=0.5in, left=0.5in, right=0.5in]{geometry}
\usepackage{enumitem}
\usepackage[utf8]{inputenc}
\usepackage[english]{babel}
\usepackage{multicol}
\usepackage{fontawesome}
\usepackage{hyperref}

\setlength{\parindent}{0pt}



\begin{document}
\begin{center}
\thispagestyle{empty}
\begin{Large}
    \textbf{Robin Uhrich\\}
\end{Large}
\begin{center}   
    \begin{tabular}{c l l c l}
         \normalsize \faMapMarker & University of Tübingen & \hfill & \faGraduationCap &  M.Sc. Machine Learning \\ 
         \normalsize \faAt & \href{mailto:robin.uhrich@gmail.com}{robin.uhrich@gmail.com} 
         & \hfil & \normalsize \faGlobe & \href{https://robinu434.github.io/}{https://robinu434.github.io/} \\ 
         % \faPhone & +49 172 9606083 \\
         \normalsize \faGithubSquare &  \href{https://github.com/RobinU434/}{RobinU434} & \hfil & \normalsize \faLinkedinSquare & \href{https://www.linkedin.com/in/robin-uhrich-bb2891293}{robin-uhrich-bb2891293}
    \end{tabular}
    \\ 
\end{center}
    \hrulefill
\end{center}


%%%%%%%%%%%%%%%%%%%%%%%%%%%%%%%%%%%%%%%%%%%%%%%%%%%%%%%%%%%%%%%
% OBJECTIVE
% Who you are, what domain, what are you looking for and when?
%%%%%%%%%%%%%%%%%%%%%%%%%%%%%%%%%%%%%%%%%%%%%%%%%%%%%%%%%%%%%%%
%\noindent \textbf{\underline{OBJECTIVE}} \\
%\noindent I am an ungraduated student in computer science in the second year at the Albert Ludwigs Universität Freiburg. I \\
% Skill 1 (level of expertise), Skill 2 (level of expertise), Skill 3 (level of expertise) \\

%%%%%%%%%%%%%%%%%%%%%%%%%%%%%%%%%%%%%%%%%%%%%%%%%%%%%%%%%%%%%%%
% EDUCATION
% University name, degree, year of graduation, GPA (optional)
%%%%%%%%%%%%%%%%%%%%%%%%%%%%%%%%%%%%%%%%%%%%%%%%%%%%%%%%%%%%%%%
\noindent \textbf{\underline{EDUCATION}} \\

% \noindent\textbf{High School Backnang} \hfill Backnang, Germany \\
% \textit{Technical High School} \hfill September 2015 $-$ July 2018 \\ \\
\noindent\textbf{University of Tübingen} \hfill Tübingen, Germany  \\
\textit{Master Machine Learning} \hfill since October 2023
\vspace{0.3cm}\\
\noindent\textbf{Politecnico di Milano} \hfill Milano, Italy  \\
\textit{Erasmus Exchange, Master Computer Science} \hfill February 2025 $-$ August 2025
\vspace{0.3cm}\\
\textbf{University of Freiburg} \hfill Freiburg im Breisgau, Germany  \\
\textit{Bachelor Computer Science} \hfill October 2019 $-$ September 2023
\vspace{0.3cm}\\
\textbf{University of Stuttgart} \hfill Stuttgart, Germany \\
\textit{Aeronautics and Space Engineering} \hfill October 2018 $-$ September 2019 
\vspace{0.2cm}\\


%%%%%%%%%%%%%%%%%%%%%%%%%%%%%%%%%%%%%%%%%%%%%%%%%%%%%%%%%%%%%%%
% PUBLICATIONS
%%%%%%%%%%%%%%%%%%%%%%%%%%%%%%%%%%%%%%%%%%%%%%%%%%%%%%%%%%%%%%%
\noindent \textbf{\underline{PUBLICATIONS}} \\

\noindent \textbf{Occupancy Modeling via Drifting: One-Step Generation Enables Direct Policy Extraction} \hfill 2026 \\
\underline{Robin Uhrich}, Marco Bagatella, Núria Armengol Urpí, Pavel Kolev, Georg Martius \\
\textit{Preprint}, \href{https://arxiv.org/abs/2610.19883}{arXiv:2610.19883}
\vspace{0.4cm}\\

%%%%%%%%%%%%%%%%%%%%%%%%%%%%%%%%%%%%%%%%%%%%%%%%%%%%%%%%%%%%%%%
% PROJECT
% What did you do?
% How did you do it? -> Skills and technologies
% What impact did you create? -> Numbers and percentages.
%
% Talk about team work, initiative, soft skills.
%
% Can also include personal projects, competitions, contribution to Open source.
%%%%%%%%%%%%%%%%%%%%%%%%%%%%%%%%%%%%%%%%%%%%%%%%%%%%%%%%%%%%%%%
\noindent \textbf{\underline{PROJECT WORK}} \\

% \noindent \textbf{Umbrella Satellite Antenna} \textit{High School Project} \hfill  October 2016 $-$ July 2017
% \begin{itemize}[noitemsep,nolistsep,leftmargin=*]
%     \item {Collaboration with TESAT Spacecom GmbH.}
%     \item {developing a high compact parabolic satellite antenna using an umbrella}
%     \item {using rapid prototyping and mechanical CAD systems\\}
% \end{itemize}


\textbf{University} \\


% \noindent \textbf{Wall$\land$E} \textit{System Design Project} \hfill  September 2019 $-$ February 2020
% \begin{itemize}[noitemsep,nolistsep,leftmargin=*]
%     \item {Course project on developing a robot using Lego EV3 set \\ 
%     to solve multi-task problem settings.}
%     \item {Programming directly on the EV3 computer in $\mu$-Python. }
%     \item Implementation of a PID controller for line following based \\ on light and ultrasonic sensors.
%     \item Design and implementation of a finite state machine for different tasks.\\
% \end{itemize}
% 
% \noindent \textbf{2048} \textit{Game development in C++ class} \hfill  July 2020 $-$ August 2020
% \begin{itemize}[noitemsep,nolistsep,leftmargin=*]
%     \item Course project within the lecture on "Advanced Programming" \\ 
%     organized by Prof. Hannah Bast
%     \item Implemented an object-oriented version of the game "2048" in C++.
%     \item Implementation of a CI/CD pipeline, including tests and documentation.\\
% \end{itemize}

\noindent \textbf{Master Thesis} \hfill since October 2025
\begin{itemize}[noitemsep,nolistsep,leftmargin=*]
    \item Supervision by Prof. Georg Martius, Dr. Pavel Kolev, \\
    Núria Armengol Urpí, and Marco Bagatella.
    \item Topics: occupancy modelling, generative models, offline meta-RL.
    \item Resulted in the preprint \href{https://arxiv.org/abs/2610.19883}{arXiv:2610.19883} (see Publications).\\
\end{itemize}

\noindent \textbf{Distributed RL on Optimal Transport} \textit{Research Project} \href{https://github.com/RobinU434/SwitchFL}{\faGithubSquare} \hfill April 2025 $-$ October 2025
\begin{itemize}[noitemsep,nolistsep,leftmargin=*]
    \item Project under Prof. Marcello Restelli and with Dr. Marco Mussi.
    \item Topics: environment design, multi-agent RL, Flatland3.\\
    % \item Project done at the lab of Prof. Marcello Restelli supervised Dr. Marco Mussi. 
    % \item Redesign of rail environment, from Flatland3, for multi-agent \\
    % RL in an asynchronous setting to solve the train-scheduling problem.
    % \item Switch perspective from train-focused to switch-focused.
    % \item Tools used: Python, Flatlands3, Gymnasium, PettingZoo.\\
\end{itemize}

\noindent \textbf{Field Robot Event} \textit{Multi-Disciplinary Project at PoliMi}\hfill March 2025 $-$ June 2025
\begin{itemize}[noitemsep,nolistsep,leftmargin=*]
    \item Team project under Prof. Matteo Matteucci and Dr. Simone Mentasti.
    \item Topics: navigation, deep visual object tracking, decision making.\\
    % \item Course project on robotics under Prof. Matteo Matteucci, \\
    % Dr. Simone Mentasti and Mirko Usuelli.
    % \item Building and programming a robot for a \\
    % agricultural-robotics competition in a team of nine.
    % \item Implementation of navigation, computer-vision \\ 
    % and decision making algorithms.
    % \item Conducted study on visual object tracking in a \\
    % moving stereo-camera setting with DeepSORT. 
    % \item Tools used: Python, ROS2, PyTorch, Yolo\\
\end{itemize}

\noindent \textbf{Multi-Agent RL in Laser Hockey} \textit{Course Project} \href{https://github.com/RobinU434/LaserHockeyProject}{\faGithubSquare} \hfill December 2024 $-$ February 2025
\begin{itemize}[noitemsep,nolistsep,leftmargin=*]
    \item Course project by Prof. Georg Martius.
    \item Topics: SAC, Dyna, replay buffers, self-play, entropy regularization.\\
    % \item Course project on "Reinforcemnet Learning" taught by Prof. Georg Martius.
    % \item Organized competition in lazer-hockey in multi-agent RL. 
    % \item Extended SAC from stable-baselines3 with replay buffer.
    % \item Implemented and extended Deep-DynaQ with entropy regularization. 
    % \item Contributed to development of lazer-hockey environment.\\ 
\end{itemize}

\noindent \textbf{Generalized Simulation Based Inference} \textit{Research Project} 
\href{https://github.com/mackelab/neuralgbi_diffusion}{\faGithubSquare} \hfill October 2024 $-$ April 2025
\begin{itemize}[noitemsep,nolistsep,leftmargin=*]
    \item ML-Lab project under Prof. Jakob Macke and with Julius Vetter.
    \item Topics: diffusion modeling, Bayesian and probabilistic learning.\\
    % \item Extended sampling strategies in simulation based inference with \\
    % amortized cost estimation presented by Dr. Richard Gao  (\hyperlink{https://arxiv.org/abs/2305.15208}{paper}) \\
    % with guided diffusion models.
    % \item Re-implemented GBI-ACE and used it for benchmarking \\ 
    % on toy and real world datasets.
    % \item Implementation was done using the PyTorch framework for deep learning.\\
\end{itemize}

\noindent \textbf{LR - Scheduler Analysis} \textit{Seminar Project} \href{https://github.com/RobinU434/LRSchedulerAnalysis}{\faGithubSquare}\hfill April 2024 $-$ July 2024
\begin{itemize}[noitemsep,nolistsep,leftmargin=*]
    \item Seminar project under Prof. Philipp Hennig and Dr. Frank Schneider.
    \item Topics: Vision Transformers, hyperparameter search, CIFAR-10.\\
    % \item Project work for the seminar: "Deep Learning Research Kitchen" \\ 
    % hosted by Prof. Philipp Hennig, Dr. Frank Schneider, \\
    % Dr. Jonas Geiping, Dr. Antionio Orvieto.
    % \item Cunducted study of multiple learning rate scheduler and their \\ 
    % hyperparameters effecting the training of a Vision-Transformer \\
    % model trained on CIFAR-10. 
    % \item Implemented pipeline using Pytorch-Lightning, Tensorboard, \\
    % Numpy and Matplotlib for training and analysis. \\ 
\end{itemize}


\noindent \textbf{Precipitation Prediction Analysis} \textit{Course Project} \href{https://github.com/RobinU434/PrecipitationPredictionAnalysis}{\faGithubSquare}\hfill December 2023 $-$ February 2024
\begin{itemize}[noitemsep,nolistsep,leftmargin=*]
    \item Course team project under Prof. Philipp Hennig and Emilia Magnani.
    \item Topics: data science, climate data, data collection.\\
    % \item Data science project on climate data in  a group of our. 
    % \item Performed steps: Problem identification, \\ data collection, and data analysis. 
    % \item Analyzed the accuracy of the precipitation forecast \\
    % of the German weather agency against real data for the \\
    % state of Baden-Württemberg, Germany.\\
\end{itemize}

\noindent \textbf{EASRL-LM-2D-IKB} \textit{Bachelor Thesis} \href{https://github.com/RobinU434/IK-RL}{\faGithubSquare \ environment} \href{https://github.com/RobinU434/SPARK/tree/v0.0.1-bachelor-thesis-submission}{\faGithubSquare \ pipeline}\hfill July 2023 $-$ September 2023
\begin{itemize}[noitemsep,nolistsep,leftmargin=*]
    \item Project under Prof. Joschka Boedecker and with Jasper Hoffmann.
    \item Topics: variational inference, reinforcement learning, SAC, inverse kinematics.\\
    % \item Topic: Expanding Action Space in Reinforcement Learning through \\ Latent Models: A 2D % Inverse Kinematics Benchmark.
    % \item  Merging Soft-Actor-Critc with generative models such as (pretrained) \\
    % Variational Autoencoder for noise augmentation and faster learning.
    % \item Implementation of a custom inverse kinematics environment \\ 
    % with a scalable action space. \\
\end{itemize} 

\noindent \textbf{Learning Quadrupedal Locomotion From Scratch} \textit{Research Project} \href{https://github.com/feinran/locomotion_simulation}{\faGithubSquare}\hfill  April 2022 $-$ September 2022
\begin{itemize}[noitemsep,nolistsep,leftmargin=*]
    \item Research project under Prof. Joschka Boedecker and with Jasper Hoffmann.
    \item Topics: reinforcement learning, PPO, robotics, sim-to-real, reward shaping.\\
    % \item Bachelor research project together with Felix Schulz \\ at the lab of Prof. Joschka Boedecker. Supervised by Jasper Hoffmann.
    % \item Implemented training pipeline for reinforcement learning \\ 
    % experiments around PPO concerned with perturbation and reward-shaping.
    % \item Minimalist re-implementation of \hyperlink{https://arxiv.org/abs/2109.14549}{\textit{Vision-Guided Quadrupedal Locomotion}}
    % \item Used and extended an already existing physics-simulator \\ 
    % which is based on PyBullet.
    % \item Main libraries: Stable-Baselines 3, PyTorch, Gym, \\ Weights-and-Biases, and Tensorboard
    % \item Conducted a study on closing the sim to real gap on the Unitree A1 robot. \\\
\end{itemize}


\noindent \textbf{ALE Plots} \textit{Seminar Project} \href{https://github.com/RobinU434/AccumulatedLocalEffectPlots}{\faGithubSquare} \hfill  May 2022 $-$ September 2022
\begin{itemize}[noitemsep,nolistsep,leftmargin=*]
    \item Seminar project under Dr. Janek Thomas.
    \item Topics: interpretable machine learning, accumulated local effect plots.\\
    % \item Project work for the seminar: "Fair and Interpretable Machine Learning" \\  hosted by .
    % \item Re-implementation of \hyperlink{https://academic.oup.com/jrsssb/article/82/4/1059/7056085?login=true}{Accumulated local effect plots}.
    % \item Implementation was done using Python alongside Numpy and Matplotlib. \\
\end{itemize}

% \noindent \textbf{Humboldt} \textit{Game development in class ``software practical training''} % \hfill  October 2020 $-$ April 2021
% \begin{itemize}[noitemsep,nolistsep,leftmargin=*]
%     \item Course team-project under Prog. Andreas Podelski and Dr. Vincent Langenfeld.
%     \item Topics: C\#, path planning, real-time-multi-agent system.  
%     % \item Designed, implemented, and tested the game real-time \\  multi-agent strategy game: % % "Humboltd".
%     % \item Collaborated in a team of six and used Git for \\  version control.
%     % \item Programmed in C\# using the XNA-(monogame)-framework. \\
% \end{itemize}



\textbf{Student Assistant} \\

\noindent \textbf{Human2Humanoid VLA} \textit{Research Project} \hfill October 2025 $-$ March 2026
\begin{itemize}[noitemsep,nolistsep,leftmargin=*]
    \item Project under Prof. Georg Martius and Prof. Wieland Brendel.
    \item Topics: tele-operation, re-targeting, RL, action representation, \\
    VLA finetuning, human-to-robot generalization.\\
\end{itemize}

\noindent \textbf{Hand Kinematics Estimation from EMG Data} \textit{Software Development} \hfill August 2024 $-$ February 2025 
\begin{itemize}[noitemsep,nolistsep,leftmargin=*]
    \item Student assistant under Dr. Cosima Prahm and with Daniel Höglinger.
    \item Topics: deep learning, EMG kinematics data, sequential modeling.\\
    % \item Design a deep learning pipeline for \\ hand kinematics estimation based on EMG data.
    % \item Done pipeline design for a non-computer science-related researcher.
    % \item Involved redesign of dataset acquisition and processing, \\ 
    % network architecture setup, and augmentation handling.
    % \item Pipeline based on pytorch-lightning. 
    % \item Using datasets like: NinaPro, SEED, and emg2pose.\\
\end{itemize}

\noindent \textbf{Vision Imitation Learning for Modular Robots} \textit{Research Project} \href{https://github.com/RobinU434/TaskParameterizedGaussianMixtureModels}{\faGithubSquare}\hfill May 2023 $-$ August 2023
\begin{itemize}[noitemsep,nolistsep,leftmargin=*]
    \item Internship under Dr. Andrea Giusti and Dr. Marco Todescato.
    \item Topics: probabilistic learning, robotics, object tracking, distributed systems.\\
    % \item Implementation of a computer vision pipeline for \\  3D hand trajectory tracking, object recognition, and \\ segmentation with multi-stereo camera feed.
    % \item Re-implemtenation of Task-Paramterized-Gaussian Mixture model based on \\ work Silvan Calinon: \hyperlink{https://calinon.ch/papers/Calinon-JIST2015.pdf}{(paper)} for learning. 
    % \item Development of Modular-Robot optimization pipeline.
    % \item Web-based interface design for containerized pipeline setup.  
    % \item Used technologies: Python, Numpy, OpenCV, ROS2, C++, Moveit, and Eigen.\\
\end{itemize}


% \noindent \textbf{PACT} \textit{Deep Energy Retrofit} \hfill August 2022 $-$ March 2023
% \begin{itemize}[noitemsep,nolistsep,leftmargin=*]
%     \item Student assistent under Prof Alexander Reiterer /& Dr. Annette Schmitt.
%     \item Collected 3D data for point cloud segmentation with laser scanners. \\
% \end{itemize}


\noindent \textbf{Automated Pavement Operation} \textit{Research Project} \hfill  September 2021 $-$ November 2022
\begin{itemize}[noitemsep,nolistsep,leftmargin=*]
    \item Student assistant under Prof. Alexander Reiterer and Dr. Annette Schmitt.
    \item Topics: deep learning, computer vision, real-time systems, HPO.\\
    % \item Project done with Fraunhofer-IPM team under Prof. Alexander Reiterer, \\
    %  with Dr. Annette Schmitt and international project-partners.
    % \item Implemented multilabel image classification pipeline on laser scan data \\
    % with convolutional neural networks (ResNet \& VGG) and Vision-Transformer.
    % \item Conducted classification-oriented real-time and classification study with \\
    % hyper-parameter-optimization (Bayesian-Optimization + Hyperband).
    % % \item Conducted a study on real-time classification capabilities for image classification.
    % \item Implemented image processing pipeline: laser scanner $\to$ geojson \\
\end{itemize} 


% \noindent \textbf{Molecular Tumor Board Research Database} \textit{Software Development Project}\hfill  March 2020 $-$ May 2021
% \begin{itemize}[noitemsep,nolistsep,leftmargin=*]
%     \item Student assistant under Prof. Melanie Börries \& Dr. Patrick Metzger.
%     \item Topics: Database design, SQL \& No-SQL, distributed systems, Rest-API design. \\
% \end{itemize}



\vspace{0.5cm}
\textbf{Private} \\

\noindent \textbf{LazySlurm} \href{https://github.com/RobinU434/LazySlurm}{\faGithubSquare} \hfill since November 2025
\begin{itemize}[noitemsep,nolistsep,leftmargin=*]
    \item Terminal UI for monitoring Slurm HPC jobs, published on PyPI.
    \item Topics: TUI design (Textual), HPC job monitoring, remote SSH backend.\\
\end{itemize}

\noindent \textbf{Simplified RNNs} \hfill since April 2025
\begin{itemize}[noitemsep,nolistsep,leftmargin=*]
    \item Simplification of the GateL0rd architecture by Dr. Christian Gumbsch.
    \item Topics: recurrent architectures, latent state similarity.\\
\end{itemize}

\noindent \textbf{SPARK} \href{https://github.com/RobinU434/SPARK}{\faGithubSquare} \hfill April 2025 $-$ November 2025
\begin{itemize}[noitemsep,nolistsep,leftmargin=*]
    \item Stochastic Policies Augmented for Robust Knowledge/Exploration.
    \item Topics: reinforcement learning, generative modelling, SAC.\\
    % \item Build on top of the idea of my Bachelor's Thesis.
    % \item Augmentation of stochastic policies for robust exploration and more expressive action % distributions.
    % \item Extension of SAC to leverage non-diagonal Gaussian distributions for action sampling.
    % \item Benchmarking done on 2D inverse kinematics environment.
    % \item Using StableBaselines3 as the basline for SAC. \\
\end{itemize} 


\noindent \textbf{Tide} \href{https://github.com/RobinU434/Tide}{\faGithubSquare}\hfill since February 2025 
\begin{itemize}[noitemsep,nolistsep,leftmargin=*]
    \item Torch Integrated Data Engine.
    \item Topics: time series data toolkit, feature engineering.\\
    % \item Data toolkit for handling sequential data \\
    % for single or multi-file datasets. 
    % \item Compatible to PyTorch dataset, data-loader integration, \\
    % feature engineering, time-based splitting, filtering, \ldots.
    % \item Integration into publicly available time series data: \\
    % EMG, climate, finance, robotics \ldots\\
\end{itemize} 


\noindent \textbf{C2C} \href{https://github.com/RobinU434/Config2Class}{\faGithubSquare}\hfill since October 2024
\begin{itemize}[noitemsep,nolistsep,leftmargin=*]
    \item Configuration files to structured Python classes.
    \item Topics: configuration management, VS Code integration, code generation.\\
    % \item Adds support for code completion in IDE.
    % \item Integrates with Hydra and OmegaConf.\\
\end{itemize}
 


\noindent \textbf{CLIfy} \href{https://github.com/RobinU434/PyArgWriter}{\faGithubSquare}\hfill since September 2023
\begin{itemize}[noitemsep,nolistsep,leftmargin=*]
    \item Command line interface forge for Python.
    \item Topics: automatic code parsing and generation, automatic argument parsing.\\
\end{itemize}

\noindent \textbf{Hexa-ROS} \hfill since May 2021
\begin{itemize}[noitemsep,nolistsep,leftmargin=*]
    \item Hexapodal locomotion in ROS2 for the Freenove Big-Hexapod robot.
    \item Topics: ROS2, Jetson Nano, stereo vision, SLAM, locomotion, electronic design.\\
    % \item Building custom Hexapod robot using the Freenove Big-Hexapod robot kit, \\
    % Jetson Nano B01, Stereo vision, and custom stationary power supply. 
    % \item Implementing RViz2 and Gazebo simulation plugins, ROS2 architecture, \\
    % low-level servo controller, high-level walking controller. 
    % \item Designing custom proprioceptive sensing for closed-loop \\
    % controller design for cheap-PWM sensors.
    % \item Conducted a study on mobile neural-depth estimation for stereo-vision. \\
\end{itemize}

%%%%%%%%%%%%%%%%%%%%%%%%%%%%%%%%%%%%%%%%%%%%%%%%%%%%%%%%%%%%%%%
% WORK EXPERIENCE
% What did you do? -> Project goals OR what problem did you solve?
% How did you do it? -> Skills and technologies
% What impact did you create? -> Numbers and percentages.
% Example: 
% + Developed an app for matching mentor and mentees for Android and iOS platform.
% + Successfully matched 85% of the applications and randomized the rest.
% 
% Talk about team work, initiative, soft skills.
%
% Can also include personal projects, competitions, contribution to Open source.
%%%%%%%%%%%%%%%%%%%%%%%%%%%%%%%%%%%%%%%%%%%%%%%%%%%%%%%%%%%%%%%
\noindent \textbf{\underline{ACADEMIC WORK EXPERIENCE}} \\
% \noindent \textbf{Metal Works Schick, metalwork} \hfill Aspach, Germany \\
% \textit{Internship} \hfill June 2018 $-$ July 2018
% \begin{itemize}[noitemsep,nolistsep,leftmargin=*]
%     \item {Collected experience in mechanical design}
%     \item {Collected experience in metal processing\\}
% \end{itemize}% 

% \noindent \textbf{Murrplastik GmbH, injection moulding} \hfill Sulzbach, Germany \\
% \textit{Internschip} \hfill July 2018
% \begin{itemize}[noitemsep,nolistsep,leftmargin=*]
%     \item {Collected experience in design of thermoplastic resin products}
%     \item {Gained knowledge about thermoplastic characteristics\\}
% \end{itemize}% 

% \noindent \textbf{Jagic Prototype Construction, } \hfill Mittelbrüden, Germany\\
% \textit{Internship} \hfill April 2019
% \begin{itemize}[noitemsep,nolistsep,leftmargin=*]
%     \item {Collected experience in  prototypes design}
%     \item {Assist in manufacturing process}
%     \item {Collected experience in detailed metal processing\\}
% \end{itemize}


\noindent \textbf{University of North Carolina at Chapel Hill} - \textit{Visiting Researcher (upcoming)} \hfill Chapel Hill, NC, USA \\
\noindent Group of Prof. Weitong Zhang \hfill November 2026 $-$ December 2026\\

\noindent \textbf{University of Tübingen} - \textit{Student Assistant} \hfill Tübingen, Germany \\
\noindent Autonomous Learning Group  - Prof. Georg Martius \hfill October 2025 $-$ March 2026\\

\noindent \textbf{BG Clinic - Tübingen}  - \textit{Student Assistant}\hfill Tübingen, Germany \\
\noindent Hand, plastic, reconstructive, and burn surgery - Dr. Cosima Prahm  \hfill August 2024 $-$ February 2025\\

% \noindent \textbf{Volatilis}  - \textit{Freelancer} \hfill Ludwigsburg, Germany \\
% \noindent Software Development \hfill July 2024 $-$ February 2025\\
% 
% \noindent \textbf{AI-Center Tübingen}  - \textit{Student Assistant} \hfill Tübingen, Germany \\
% \noindent Central Office \hfill December 2023 $-$ May 2024 \\

\noindent \textbf{Fraunhofer Italia} - \textit{Internship} \hfill Bolzano, Italy \\
\noindent Advanced Robotics  - Dr. Andrea Giusti \hfill April 2023 $-$ August 2023 \\

\noindent \textbf{Intelligent Machine-Brain Interfacing Technology} - \textit{Student Assistant} \hfill Freiburg im Breisgau, Germany \\
\noindent Neurorobotics  - Prof. Joschka Boedecker \hfill July 2022 $-$ April 2023\\


\noindent \textbf{Fraunhofer Institute for Physical Measurement Technology} - \textit{Student Assistant}  \hfill Freiburg im Breisgau, Germany \\
\noindent Object and Shape Detection  - Prof. Alexander Reiterer \hfill June 2021 $-$ March 2023\\

\noindent \textbf{University Clinic Freiburg im Breisgau}  - \textit{Student Assistant} \hfill Freiburg im Breisgau, Germany \\
\noindent Institute of Medical Bioinformatics and System Biology - Prof. Melanie Börries \hfill March 2020 $-$ May 2021 \\ \\

%%%%%%%%%%%%%%%%%%%%%%%%%%%%%%%%%%%%%%%%%%%%%%%%%%%%%%%%%%%%%%%
% SKILLS: Important and relevant to the job you are applying for
%%%%%%%%%%%%%%%%%%%%%%%%%%%%%%%%%%%%%%%%%%%%%%%%%%%%%%%%%%%%%%%
% \vspace{0.4cm}
\noindent \textbf{\underline{TECHNICAL SKILLS}}
\begin{multicols}{2}
\begin{itemize}
    \item Python (NumPy, PyTorch (Lightning), Gymnasium, SciKit-Learn, Pandas, SB3, PettingZoo, \ldots)
    \item C, C++, C\#
    \item (NO-)SQL-Databases
    \item Docker
    \item ROS, ROS2
    \item Git
    \item SVN
    \item Mechanical CAD (Autodesk Inventor)
    \item Embedded systems (Arduino, ESP)
    \item SCRUM
    \item SLURM
\end{itemize}
\end{multicols}
\vspace{0.4cm}

%Programming in Python (one and a half years of experience), Prog (years of experience), Skill 3 (years of experience) \\

\noindent \textbf{\underline{AWARDS}}

\begin{itemize}[noitemsep,nolistsep,leftmargin=*]
\item \textbf{Field Robot Event} \textit{3rd place} \hfill August 2025
\item \textbf{System Design Project} \textit{3rd place}\hfill February 2020\\
\end{itemize}

%%%%%%%%%%%%%%%%%%%%%%%%%%%%%%%%%%%%%%%%%%%%%%%%%%%%%%%%%%%%%%%
% Extra Curricular Activities, Leadership, etc 
%%%%%%%%%%%%%%%%%%%%%%%%%%%%%%%%%%%%%%%%%%%%%%%%%%%%%%%%%%%%%%%
\noindent \textbf{\underline{EXTRA CURRICULAR ACTIVITIES}}

\begin{itemize}[noitemsep,nolistsep,leftmargin=*]
\item \textbf{Private Statistics Tutor} \hfill September 2025 $-$ October 2025
\item \textbf{Volunteer at the European Workshop on Reinforcement Learning 18} \hfill September 2025
\item \textbf{Organizer Machine Learning Reading Club} \hfill February 2025 $-$ June 2025 \\
\end{itemize}



% \noindent \textbf{Activity/ Role} \hfill Month, Year $-$ Month, Year
% \begin{itemize}[noitemsep,nolistsep,leftmargin=*]
% \item {What did you do, how did you do it and what did you achieve? \\}
% \end{itemize}

%%%%%%%%%%%%%%%%%%%%%%%%%%%%%%%%%%%%%%%%%%%%%%%%%%%%%%%%%%%%%%%
% Other Skills: you can add all your other skills here.
% Continue to keep only relevant skills
%%%%%%%%%%%%%%%%%%%%%%%%%%%%%%%%%%%%%%%%%%%%%%%%%%%%%%%%%%%%%%%
\noindent \textbf{\underline{LANGUAGES}} \\ 

German (native), English (C1/C2)
\vspace{0.25cm} \\

\noindent \textbf{\underline{PERSONAL INTERESTS}} \\

Archery (bronze medal at European Championships 2018), Running, Road cycling, Photography
% \begin{multicols}{1}
% \begin{itemize}
%     \item Archery (bronze medal at European Championships 2018)
%     \item Running (10 km -20 km)
%     \item Road cycling
%     \item Photography
% \end{itemize}
% \end{multicols}


\end{document}

